\documentclass{ltjsarticle}
\begin{document}
\title{圧縮性CFDのためのCUDA Fortran勉強会資料}
\author{岡山大学　空気力学研究室　幡山　純}
\maketitle
\clearpage

\tableofcontents
\clearpage

\part{CUDA Fortran概要}
\section{GPUの計算ユニット}
\input{./part1/gpu_compute_unit.tex}

\section{GPUのメモリ階層構造}
\input{./part1/gpu_memory.tex}

\clearpage

\part{OUxSBLI解説}
\section{圧縮性流体の数値計算法}
\subsection{支配方程式}
\input{./part2/equations.tex}

\subsection{空間離散化}
\input{./part2/spatial.tex}

\subsection{時間離散化}
\input{./part2/temporal.tex}

\section{OUxSBLIの計算フロー}
\input{./part2/flow.tex}

\section{GPUへのメモリコピー}
CUDA CでGPU上に配列を確保するにはcudaMallocを使用する．CUDAは一次元配列の扱いが得意であるため，多次元配列を扱う際でもメモリ上で一次元に平坦化されたメモリを思い浮かべながらプログラムを書くとよい．CUDA CではcudaMallocでGPU上に場所を確保した後にcudaMemcpyでCPUからGPUへのコピーを行う．CUDA CではcudaMallocとcudaMemcpyの両方を明示的に書く必要があるが，CUDA Fortranでは以下のようにdevice属性を付与して配列を宣言すればGPU上に配列が確保される．（内部的にはcudaMallocが走る）

\begin{lstlisting}[language=Fortran]
real(8), device :: Q_gpu(5,nx,ny,nz)
\end{lstlisting}

\noindent
配列を宣言したうえで通常の代入文を書けばcudaMemcpyを実行できる．なお，GPU上での配列の宣言はallocatable属性を付与して動的に行うこともできる．

CUDA Fortranではconstant memoryを活用することが重要である．使い方は簡単で，parameter属性かconstant属性を付けておけばconstant memory上に確保される．Constant memoryは複数のスレッドが同時にアクセスしてもbank conflict(後述)を起こさないという利点をもつ．また，GPUは割り算が非常に遅いため頻繁に使う分数の逆数を事前に計算しておくのも大切な高速化手法であるといえる．

\begin{lstlisting}[language=Fortran]
real(8), parameter :: one_third = 1.d0 / 3.d0
real(8), parameter :: two_third = 2.d0 / 3.d0
real(8), parameter :: one_sixth = 1.d0 / 6.d0
\end{lstlisting}


\section{保存量から基本量を計算～directiveベースの並列化～}
CUDA FortranでもOpen ACCのようなdirectiveベースの並列化を実装できる．コンパイラの判断でimplicit private等の属性を付与してくれる．Reduction処理を扱うこともできる．コンパイラオプションに-Minfo=accelを付与することでコンパイラがどのように並列化したかを確認することができる．

OUxSBLIでは保存量から基本量を計算する処理にdirectiveベースの並列化を用いている．以下にコードを示す．

\begin{lstlisting}[language=Fortran]
!$cuf kernel do(3)<<<*,(32,4,2)>>>
do k = 1, nz
  do j = 1, ny
    do i = 1, nx
      over_Q1 = 1.d0 / QJ(1,i,j,k)
      rho     = Jacobian(i,j) * QJ(1,i,j,k)
      u       = QJ(2,i,j,k) * over_Q1
      v       = QJ(2,i,j,k) * over_Q1
      w       = QJ(2,i,j,k) * over_Q1
      p       = gamma_1 * (Jacobian(i,j) * QJ(5,i,j,k) - 0.5d0 * rho * (u*u + v*v + w*w))
      Q(1,i,j,k) = rho
      Q(2,i,j,k) = u
      Q(3,i,j,k) = v
      Q(4,i,j,k) = w
      Q(5,i,j,k) = p
      T(i,j,k)   = p / (R * rho)
enddo;enddo;enddo
\end{lstlisting}

\noindent
Directiveベースの並列化ではコンパイラは”安全”な並列化を行うため，本質的にカーネルを起動する並列化と同じコードを書いたとしてもレジスタの使用量等に違いがでる場合がある．また，GPUは割り算がとにかく苦手なので，事前に逆数を計算するのがよい．


\section{数値流束の計算～カーネルを使った並列化～}
\subsection{カーネルの使い方}
CUDA Fortranでは以下のように記述するとカーネルが呼び出される．

\begin{lstlisting}[language=Fortran]
call calc_keep_x<<<blocks, threads, stream>>>(nx, ny, nz, ruvwp, T, E)
\end{lstlisting}

三重の不等号はシェブロンと呼ばれる．カーネルを呼ぶにはブロック数，スレッド数，実行するストリームを指定する．デフォルトのストリームは0番である．ストリーム番号を変えるとカーネルが並行実行されることがある（計算リソースが不足するため，実際に並行実行されることはほとんどない）．

blocksとthreadsは以下のように宣言する．

\begin{lstlisting}[language=Fortran]
type(dim3), parameter :: threadsE = dim3(32, 2, )
type(dim3), parameter :: blocksE  = dim3((nx-1+threadsE%x-1)/threadsE%x, &
                                         (ny-1+threadsE%y-1)/threadsE%y, &
                                         (nz-1+threadsE%z-1)/threadsE%z)
\end{lstlisting}

\noindent
シェブロンで呼び出されるsubroutineにはglobal属性を付与する．

\begin{lstlisting}[language=Fortran]
attributes(global) subrouinte calc_keep_x(nx, ny, nz, ruvwp, T, E)
\end{lstlisting}

\noindent
global属性が付与されたsubroutine内ではblockIdx, blockDim, threadIdxが使用できる．

\begin{itemize}
\item blockIdx:現在実行されているblockの番号
\item blockDim:blockに含まれるthreadの数
\item threadIdx:block内でのthread番号（最大値はblockDim）
\end{itemize}

\noindent
これらをi, j, kと対応させるには以下のように書く．

\begin{lstlisting}[language=Fortran]
i = (blockIdx%x-1)*blockDim%x + threadIdx%x
j = (blockIdx%y-1)*blockDim%y + threadIdx%y
k = (blockIdx%z-1)*blockDim%z + threadIdx%z
\end{lstlisting}


\subsection{高次精度差分におけるshared memoryの活用}
CUDA Fortranで高次精度の差分を計算する際はshared memoryの使用が推奨される．Shared memoryを使う場合はshared属性をつけて配列を宣言する．Shared memoryの形状にはparameter属性を付与してコンパイル時に形状が決まるようにするとよい．Shared memoryを動的に確保することもできるが，静的に確保した方が速い．

6点stencilの差分計算において，global memoryからshared memoryに素朴にコピーしてみる．Qをglobal memory上の配列，rhoをshared memory上の配列とする．

\begin{lstlisting}[language=Fortran]
rho(it-2:it+3,jt,kt) = Q(1,i-2:i+3,j,k)
\end{lstlisting}

\noindent
このようにすると複数のスレッドが同一のshared memory上の番地にアクセスするため”待ち”が発生する．これをBank conflictという．Bank conflictを解消するには以下のように書くとよい．なお，syncthreadsを呼び忘れるとスレッド間で同期がとられないため以降の差分計算が失敗する．

\begin{lstlisting}[language=Fortran]
it = threadIdx%x
i_base = (blockIdx%x-1)*blockDim%x
do ii = it-2, threadsE%x+3, blockDim%x
  i = i_base + ii
  if (i >= 1 .and. i <= nx) then
    rho(ii,jt,kt) = Q(1,i,j,k)
  endif
enddo
call syncthreads()
\end{lstlisting}

\noindent
CUDA Fortranで究極的な実行速度を達成するには一次元配列を用いるのが望ましい．以下に一次元のshared memoryを用いた例を示す．OUxSBLIでは以下のようにglobal memoryからshared memoryへのコピーを行っている．（2026年3月時点）

\begin{lstlisting}[language=Fortran]
it = threadIdx%x
jt = threadIdx%y
kt = threadIdx%z
j  = (blockIdx%y-1)*blockDim%y + jt + 1
k  = (blockIdx%z-1)*blockDim%z + kt + 1
i_base = (blockIdx%x-1)*blockDim%x
offset_yz = (jt-1) * sx + (kt-1) * sx * sy
do ii = it-2, threadsE%x+3, blockDim%x
  i = i_base + ii
  if (i >= 1 .and. i <= nx .and. j >= 1 .and. j <= ny .and. k >= 1 .and. k <= nz) then
    idx = ii + offset_yz
    rho(idx) = Q(1,i,j,k)
      u(idx) = Q(2,i,j,k)
      v(idx) = Q(3,i,j,k)
      w(idx) = Q(4,i,j,k)
      p(idx) = Q(5,i,j,k)
    tmp(idx) =   T(i,j,k)
  endif
enddo
call syncthreads()
\end{lstlisting}

CUDA Fortranでは連続アクセスでない部分配列をsubroutineに渡すことができない．したがって，y方向＆z方向の差分を計算する際に問題が生じる．OUxSBLIではy方向＆z方向の差分を計算する際はshared memory上で差分方向に連続アクセスになるように転置するという方式を採用している．単純な中心差分では転置のコストで全体が遅くなる可能性が考えられるが，高解像度風上差分を用いる場合は転置のコストよりも連続アクセスであることによるその後の処理の恩恵の方が大きいと考えている．

\begin{lstlisting}[language=Fortran]
it = threadIdx%x
jt = threadIdx%y
kt = threadIdx%z
i  = (blockIdx%x-1)*blockDim%x + it + 1
k  = (blockIdx%z-1)*blockDim%z + kt + 1
j_base = (blockIdx%y-1)*blockDim%y
offset_xz = (it-1) * sy + (kt-1) * sy * sx
do jj = jt-2, threadsF%y+3, blockDim%y
  j = j_base + jj
  if (i >= 1 .and. i <= nx .and. j >= 1 .and. j <= ny .and. k >= 1 .and. k <= nz) then
    idx = jj + offset_xz
    rho(idx) = Q(1,i,j,k)
      u(idx) = Q(2,i,j,k)
      v(idx) = Q(3,i,j,k)
      w(idx) = Q(4,i,j,k)
      p(idx) = Q(5,i,j,k)
    tmp(idx) =   T(i,j,k)
  endif
enddo
call syncthreads()
\end{lstlisting}


\section{時間発展～WarpShuffleを使ってみる～}
CUDA Fortranではwarp単位で計算が実行される．1つのwarpは32スレッドからなる．こうした事情からCUDA Fortranではブロックあたりのスレッド数を32の倍数にするのが良いとされている．Warp shuffle命令を使用すると同一warp内の隣のスレッドが担当するメモリにアクセスすることができる．Shared memoryを用いても隣のスレッドが担当するメモリにアクセスできるが，shared memoryへのコピーに時間を要するという問題がある．一方で，warp shuffle命令ではゼロコピーで隣のメモリを見に行くため高速である．しかし，warp shuffleでアクセスできるのはGPU上でも隣にあるメモリだけである．したがって，warp shuffleで差分が計算できるのはx方向に限定される．また，warp shuffle命令ではレジスタを消費する．したがって，レジスタ使用量が少ないRunge-Kuttaの時間発展におけるx方向の差分にwarp shuffle命令を用いて高速化を試みる（Runge-Kuttaの時間発展はボトルネックになっていないため，わざわざwarp shuffleを用いる必要は薄い．比較的簡単な処理でwarp shuffleを説明するためにここで取り上げている）．

まず，warp shuffleを使用するには

\begin{lstlisting}[language=Fortran]
use cooperative_groups
\end{lstlisting}

\noindent
が必要である．cooperative\_groupsには他にも便利な機能が含まれている（詳しくはNvidia公式のユーザーガイドを参照）．Warp shuffle命令を用いて差分計算を行っている箇所を以下に示す．

\begin{lstlisting}[language=Fortran]
do l = 1, 5
  E_curr   = E(l,i,j,k)
  tmp_bits = __shfl_down_sync(z'ffffffff', transfer(E_curr, 0_8), 1)
  E_next   = transfer(tmp_bits, 0.0_8)
  if (lane == 31 .or. i == nx-2) then
    E_next = E(l,i+1,j,k)
  endif
  R = dtdydz * (-E_curr     + E_next) &
  & + dtdzdx * (-F(l,i,j,k) + F(l,i,j+1,k)) &
  & + dtdxdy * (-G(l,i,j,k) + G(l,i,j,k+1))
  Q2(l,i+1,j+1,k+1) = Q(l,i+1,j+1,k+1) - coef1 * R
  Rs(l,i,j,k) = Rs(l,i,j,k) + coef2 * R
enddo
\end{lstlisting}

\noindent
transferを用いて第一引数のビット表現を取得する．倍精度実数型を扱っているため第二引数には0\_8を渡す．transferを使わずにwarp shuffle命令で倍精度実数型を扱うと結果が変わる恐れがある（CUDA Fortranの仕様かも）．
Warp shuffle命令で隣の値を読みに行くには\_shfl\_down\_syncを使う．1つ後ろのメモリを読みに行くため\_shfl\_donw\_syncを使う．1つ前のメモリを読みに行く場合\_shfl\_up\_syncを使う．第一引数にz'ffffffff'を指定することでwarp内で同期をとることができる．


\clearpage

\part{雑多なトピックス}
\section{マルチGPUへの対応}
\input{./part3/mpi.tex}

\section{コンパイルオプション}
\input{./part3/options.tex}

\section{プロファイラーの使用}
\input{./part3/nsys.tex}

\section{デバッガ―の使用}
\input{./part3/sani.tex}

\end{document}
